\chapter{Podsumowanie}
\label{chap:podsumowanie}
W pracy opracowano i opisano kompletną architekturę Shadow Deploy dla funkcji SMF w środowisku 5G Core. Przyjęty model łączy Kubernetes, Istio i Envoy z metryczną walidacją jakości opartą o Prometheus i Iter8.

\section{Stopień realizacji celu}
Cel pracy został zrealizowany. Przygotowane rozwiązanie spełnia założenia:
\begin{itemize}
    \item realizuje bezinwazyjne mirrorowanie ruchu do wersji kandydackiej,
    \item izoluje stan danych wersji shadow od danych produkcyjnych,
    \item dostarcza automatyczny mechanizm oceny jakości na podstawie SLO.
\end{itemize}

\section{Najważniejsze wnioski}
\begin{enumerate}
    \item Shadow Deploy w Service Mesh pozwala testować nową wersję funkcji na ruchu produkcyjnym bez ryzyka wpływu na użytkowników.
    \item Separacja bazy \texttt{udr-shadow} jest warunkiem koniecznym bezpiecznej walidacji funkcji kontrolnych 5G Core.
    \item Metryki z Envoy i Prometheus są wystarczające do budowy obiektywnego \emph{quality gate} wspierającego decyzję o promocji wersji.
\end{enumerate}

\section{Ograniczenia}
Praca koncentruje się na funkcji SMF i scenariuszach krótkookresowych. Nie obejmuje pełnego spektrum awarii wieloregionowych, aspektów kosztowych dużych klastrów operatorskich ani długoterminowej degradacji parametrów.

\section{Kierunki dalszych prac}
W dalszym etapie warto rozszerzyć rozwiązanie o:
\begin{itemize}
    \item adaptacyjne sterowanie \texttt{mirrorPercentage} zależnie od bieżącej jakości usług,
    \item automatyczne rollbacki uruchamiane przez sygnały SLO,
    \item rozszerzenie modelu na kolejne funkcje 5G Core (AMF, PCF, NRF),
    \item integrację z politykami bezpieczeństwa typu zero trust w warstwie mesh.
\end{itemize}

Przedstawione podejście stanowi praktyczny fundament do budowy bezpiecznego procesu wdrożeń CNF w środowiskach telekomunikacyjnych klasy operatorskiej.
